Bei konstantem Druck ist das Volumen einer bestimmten Gasmenge proportional zur absoluten Temperatur (Gay-Lussac): $V_2 = V_1\cdot T_2/T_1$. Die Masse bleibt gleich, und mit $\rho = m/V$ folgt
\begin{align}
 \frac{\rho_2}{\rho_1} &= \frac{V_1}{V_2} = \rqF = \frac{\TEi}{\TZw} = \rq
\end{align}
mit $T_1 = \TEiF = \TEi$ und $T_2 = \TZwF = \TZw$. Die Dichte ändert sich also um
\begin{align}
 \frac{\Delta\rho}{\rho_1} &= \drrF = \drr \approx \result{\drrP{0}{2}}
\end{align}
